%% main.tex
%% 中国农业大学学术学位博士学位论文 - 主文件
%%
%% 编译方式（推荐命令行或使用latexmk）：
%%   xelatex main
%%   bibtex main        （如使用biblatex则用 biber main）
%%   makeindex main.nlo -s nomencl.ist -o main.nls  （如有符号表）
%%   xelatex main
%%   xelatex main
%%
%% 选项说明：
%%   默认目录显示到二级标题（1.1）
%%   如需显示三级标题（1.1.1），请将下面的 [] 改为 [tocthree]

%% Author: ZHuo Yue
%% Email: zhuo0000yue@gmail.com
%% GitHub: github.com/Zhuoy3/cau-thesis-doctor
%% License: GPL-3.0 license
%% Last Modified: 2026/02/26

\documentclass[]{CAUthesis}

\addbibresource{references.bib}

%% ============================================================
%% 论文信息填写
%% ============================================================
\classno{}                                  % 分类号（选填）
\secretlevel{}                              % 密级（公开则留空）
\studentid{2021XXXXX}                       % 学号
\titlecn{论文中文题目}                       % 中文题目
\titleen{English Title of the Thesis}        % 英文题目
\author{张三}                               % 研究生姓名
\supervisor{李四 教授}                       % 指导教师
\cosupervisor{}                              % 合作指导教师（无则留空）
\degreetype{农学博士}                        % 申请学位门类级别
\major{作物栽培学与耕作学}                    % 专业名称
\research{作物高产高效栽培理论与技术}          % 研究方向
\college{农学院}                             % 所在学院
\date{2025年6月}                             % 论文完成日期

%% ============================================================
%% 符号表定义（如有需要）
%% ============================================================
% 在此定义符号，编译时需运行 makeindex
% 示例：
\nomenclature{$\alpha$}{学习率}
\nomenclature{$\beta$}{衰减系数}
\nomenclature{DNA}{脱氧核糖核酸 (Deoxyribonucleic Acid)}

%% ============================================================
%% 正文开始
%% ============================================================
\begin{document}

%% ---- 封面 ----
\makecover{}

%% ---- 独创性声明 ----
\makedeclaration{}

%% ---- 前置部分开始 ----
\frontmatter{}

%% ---- 中文摘要 ----
\begin{abstractcn}
  在此撰写中文摘要内容。

  论文摘要内容包括研究工作目的、研究方法、所取得的结果和结论；应突出本论文的创造性成果或新见解，语言精练。博士论文摘要一般不超过1500字。
\end{abstractcn}
\keywordscn{关键词一，关键词二，关键词三，关键词四，关键词五}

%% ---- 英文摘要 ----
\begin{abstracten}
  Write English abstract here.

  The abstract should include the purpose of the research, methods, results and conclusions. It should highlight the creative achievements or new insights of this thesis.
\end{abstracten}
\keywordsen{Keyword 1, Keyword 2, Keyword 3, Keyword 4, Keyword 5}

%% ---- 目录 ----
\tableofcontents{}

%% ---- 插图和附表清单（图表较多时使用） ----
\makelistoffiguresandtables{}

%% ---- 主要符号表（如有需要） ----
% 取消以下注释以启用符号表，编译时需运行 makeindex
\makesymboltable{}

%% ---- 正文部分开始 ----
\mainmatter{}

%% ============================================================
%% 第一章 引言
%% ============================================================
\chapter{引言}

\section{研究背景与意义}

在此撰写研究背景与意义的内容。正文采用小四号宋体，行间距为20磅。文中英文、罗马字符一般采用Times New Roman正体。

引用示例：产业融合\parencite{tiancaihong2024}、经济增长\parencite{zhang2023}。

\subsection{研究背景}

这是三级标题的示例内容。

\section{国内外研究进展}

在此撰写文献综述内容。

\section{研究目标与内容}

在此撰写研究目标与内容。

%% ============================================================
%% 第二章 材料与方法
%% ============================================================
\chapter{材料与方法}

\section{实验材料}

正文内容示例。

\subsection{供试材料}

更详细的材料描述。

\section{实验方法}

实验方法描述。

\section{数据分析}

数据分析方法描述。

% ---- 公式示例 ----
公式编号按章节编号，居中对齐，编号右对齐：

\begin{equation}
  E = mc^2
  \label{eq:einstein}
\end{equation}

\begin{equation}
  F(x) = \int_{-\infty}^{x} f(t) \, dt
  \label{eq:integral}
\end{equation}

由\cref{eq:einstein}可知，质量与能量……

%% ============================================================
%% 第三章 结果与分析
%% ============================================================
\chapter{结果与分析}

\section{实验结果一}

\subsection{详细结果}

% ---- 图片示例 ----
% 图形不能跨页显示，图序、标题、图例说明居中置于图的下方
\begin{figure}[htbp]
  \centering
  % \includegraphics[width=0.8\textwidth]{figures/example.png}
  \rule{0.8\textwidth}{5cm}  % 占位符，请替换为实际图片
  \caption{示例图片标题}\label{fig:example}
\end{figure}

如\cref{fig:example} 所示，实验结果表明……

% ---- 三线表示例 ----
% 三线表：顶线和底线为粗线(1.5磅)，栏目线为细线(0.75磅)
\begin{table}[htbp]
  \centering
  \caption{实验数据分析}\label{tab:example}
  \zihao{5}\songti{}  % 表格内容五号宋体
  \begin{tabular}{lccc}
    \toprule
        & 食品安全类文档数 & 非食品安全类文档数 & 总数    \\
    \midrule
    训练集 & 1898     & 11388     & 13286 \\
    验证集 & 250      & 250       & 500   \\
    测试集 & 250      & 250       & 500   \\
    \midrule
    总数  & 2398     & 11888     & 14286 \\
    \bottomrule
  \end{tabular}

  \vspace{3pt}
  {\zihao{-5}\songti{} 注：表注内容采用小五号宋体。}  % 表注：小五号宋体
\end{table}

如\cref{tab:example} 所示，实验数据……

\section{实验结果二}

\subsection{详细结果}

分析讨论内容。

\section{实验结果三}

结果描述。

%% ============================================================
%% 第四章 讨论
%% ============================================================
\chapter{讨论}

\section{主要发现}

讨论内容。

\section{与已有研究的比较}

比较分析。

\section{研究局限与展望}

局限性和未来展望。

%% ============================================================
%% 第五章 结论
%% ============================================================
\chapter{结论}

\section{主要结论}

结论应明确、简练、完整、准确。

\section{创新点}

创新点描述。

\section{展望}

未来研究展望。

%% ============================================================
%% 后置部分
%% ============================================================
\backmatter{}

%% ---- 参考文献 ----
% 参考文献一律放在结论/讨论之后
% 使用BibTeX管理参考文献
\phantomsection{}\printbibliography[heading=bibintoc, title={参考文献}]

%% ---- 致谢 ----
\begin{acknowledgements}
  在此撰写致谢内容。致谢对象限于对课题研究、学位论文完成等方面有较重要帮助的人员。

  \hfill 2025年6月
\end{acknowledgements}

%% ---- 附录 ----
\appendix

\chapter{补充数据}

附录内容。凡不宜放在论文正文中，但又与论文有关的研究过程或资料。

\chapter{程序代码}

更多附录内容。每一附录均另页起。

%% ---- 作者简介 ----
\begin{authorinfo}
  \textbf{基本介绍：}张三，男，1995年1月出生，XX省XX市人。

  \vspace{1ex}
  \textbf{教育经历：}
  \begin{itemize}[leftmargin=2em, nosep]
    \item 2013年9月—2017年6月，XX大学，XX专业，本科/学士
    \item 2017年9月—2019年6月，XX大学，XX专业，硕士
    \item 2019年9月—2025年6月，中国农业大学，XX专业，博士
  \end{itemize}

  \vspace{1ex}
  \textbf{研究生期间发表的学术论文：}
  \begin{enumerate}[leftmargin=2em, nosep]
    \item 张三，李四. 论文题目[J]. 期刊名称，2024，XX (X)：XX-XX.\
    \item Zhang S, Li S. Paper title[J]. Journal Name, 2024, XX (X): XX-XX.\
  \end{enumerate}

  \vspace{1ex}
  \textbf{研究生期间主持/参与的科研项目：}
  \begin{enumerate}[leftmargin=2em, nosep]
    \item 项目名称，项目来源，项目编号，起止时间，参与
  \end{enumerate}

  \vspace{1ex}
  \textbf{研究生期间获得的奖励和荣誉：}
  \begin{enumerate}[leftmargin=2em, nosep]
    \item 奖励名称，授予单位，年份
  \end{enumerate}

  \vspace{1ex}
  \textbf{其他成果：}

  无。
\end{authorinfo}

\end{document}
